% ===================== Topic 1 -- Algebra and Equations -- part c (EN): Problems 31--40 =====================
\providecommand{\rvabBox}[4]{\begin{tikzpicture}[baseline=-0.7ex,font=\small]
  \path (-0.66,-0.42) rectangle (0.66,0.42);
  \draw[Navy,line width=0.8pt,fill=tmBlue!6] (-0.66,-0.32) rectangle (0.66,0.32);
  \draw[Navy,line width=0.8pt] (0,-0.32)--(0,0.32) (-0.66,0)--(0.66,0);
  \node[inner sep=0pt] at (-0.33,0.16){$#1$};\node[inner sep=0pt] at (0.33,0.16){$#2$};
  \node[inner sep=0pt] at (-0.33,-0.16){$#3$};\node[inner sep=0pt] at (0.33,-0.16){$#4$};\end{tikzpicture}}

\clearpage
\dividerbar{Extension and challenge}

% ---------------------------------------------------------------- Problem 31
\begin{probbox}{31}{\relax}
The road from Mr Henry's home to the city is $180$ km long. He rides a motorbike at a constant speed, and after $3$ hours he has covered
$\dfrac{2}{3}$ of the distance.
\begin{enumerate}[label=\alph*),leftmargin=1.6em,itemsep=1pt,topsep=2pt]
\item Find the speed of the motorbike.
\item If he keeps the same speed, how much longer will it take him to reach the city?
\item If he wants to arrive $30$ minutes earlier than planned in part b), at what speed must he ride for the rest of the way?
\end{enumerate}
\end{probbox}
\probsol{
\textbf{Understanding the problem.} \emph{Given:} the whole distance $180$ km; in $3$ hours he covered $\dfrac{2}{3}$ of it.
\emph{Find:} the speed; the remaining time; the new speed needed to save $30$ minutes.
\begin{recallbox}$s=v\cdot t$;\quad $v=\dfrac{s}{t}$;\quad $t=\dfrac{s}{v}$. \ \ $\dfrac{m}{n}$ of $A$ equals $A\div n\cdot m$.\end{recallbox}
\step{1}{Distance covered and distance left}
Covered: $\dfrac{2}{3}$ of $180$ km is $180\div 3\cdot 2=60\cdot 2=120$ (km). Left: $180-120=60$ (km).
\step{2}{Part a: choose an unknown and write an equation}
Let the speed of the motorbike be $v$ (km/h), $v>0$. In $3$ hours he covers $3v$ (km). By Step 1 this distance is $120$ km, so $3v=120$, which gives $v=120\div 3=40$ (km/h).
\step{3}{Part b: the time left}
Let $t$ (hours) be the time needed for the remaining $60$ km at $40$ km/h. Then $40t=60$, so $t=60\div 40=1.5$ (hours) $=1$ hour $30$ minutes.
\begin{tipbox}Another way, using direct proportion: at a constant speed, time is directly proportional to distance. $120$ km takes $3$ hours; $60$ km is half of $120$ km, so it takes half of $3$ hours, that is $1.5$ hours.\end{tipbox}
\step{4}{Part c: the new speed}
New time: $1$ hour $30$ minutes $-\,30$ minutes $=1$ hour. Let the new speed be $u$ (km/h). Then $u\cdot 1=60$, so $u=60$ (km/h).
\begin{warnbox}Do not just add to the speed, as in ``$30$ minutes earlier means $30$ km/h faster''. Recalculate: speed $=$ distance $\div$ new time.\end{warnbox}
\verify{$40\cdot 3=120$ km; $40\cdot 1.5=60$ km; $60\cdot 1=60$ km. Total $120+60=180$ km.}
\ans{a) $40$ km/h;\ \ b) $1$ hour $30$ minutes;\ \ c) $60$ km/h.}
}

% ---------------------------------------------------------------- Problem 32
\begin{probbox}{32}{\relax}
An isosceles triangle has a base of $10$ cm and two equal sides of $2x$ cm each. A rectangle has length $(x+7)$ cm and width $5$ cm.
The perimeter of the triangle equals the perimeter of the rectangle.
\begin{figC}
\draw[tmBlue,line width=1.1pt,fill=tmBlue!10] (0,0)--(2,0)--(1,2.615)--cycle;
\draw[tmBlue,line width=0.9pt] (0.388,1.351)--(0.612,1.265) (1.388,1.265)--(1.612,1.351);
\node[below] at (1,0){$10$ cm};
\node[anchor=east] at (0.3,1.45){$2x$ cm};\node[anchor=west] at (1.7,1.45){$2x$ cm};
\draw[tmOrange!85!black,line width=1.1pt,fill=tmOrange!12] (3.6,0) rectangle (6.4,1);
\node[below] at (5,0){$(x+7)$ cm}; \node[right] at (6.4,0.5){$5$ cm};
\end{figC}
\begin{enumerate}[label=\alph*),leftmargin=1.6em,itemsep=1pt,topsep=2pt]
\item Write an expression in $x$ for the perimeter of each shape.
\item Find $x$, then find the side lengths of the triangle and the area of the rectangle.
\end{enumerate}
\end{probbox}
\probsol{
\textbf{Understanding the problem.} The triangle has sides $2x$, $2x$, $10$ (cm). The rectangle measures $x+7$ by $5$ (cm). The two perimeters are equal.
\begin{recallbox}Perimeter of a triangle $=$ sum of the three sides. Perimeter of a rectangle $=2\times(\text{length}+\text{width})$. Area of a rectangle $=$ length $\times$ width.\end{recallbox}
\step{1}{Part a: expressions for the perimeters}
Triangle: $2x+2x+10=4x+10$ (cm).\quad Rectangle: $2\cdot\big((x+7)+5\big)=2(x+12)=2x+24$ (cm).
\begin{warnbox}When multiplying the bracket by $2$, multiply both terms: $2(x+12)=2x+24$, not $2x+12$.\end{warnbox}
\step{2}{Part b: write an equation}
$4x+10=2x+24$.
\step{3}{Solve the equation}
Subtract $2x$ from both sides: $2x+10=24$. Subtract $10$ from both sides: $2x=14$. Divide by $2$: $x=7$.
\step{4}{Work out what is asked}
Equal sides of the triangle: $2x=14$ cm, so the triangle has sides $14$ cm, $14$ cm, $10$ cm.\\
Rectangle: length $x+7=14$ cm, width $5$ cm, area $14\cdot 5=70$ (cm$^2$).
\verify{Perimeter of the triangle $14+14+10=38$ cm; perimeter of the rectangle $2(14+5)=38$ cm. Equal.}
\ans{a) $4x+10$ and $2x+24$;\ \ b) $x=7$; triangle $14$ cm, $14$ cm, $10$ cm; area of the rectangle $70$ cm$^2$.}
}

% ---------------------------------------------------------------- Problem 33
\begin{probbox}{33}{\stars{1}}
We use the rule of Problem 14 again:\quad \rvabBox{a}{b}{c}{d} $=a\cdot d+b\cdot c$.
\begin{enumerate}[label=\alph*),leftmargin=1.6em,itemsep=2pt,topsep=2pt]
\item Evaluate \rvabBox{-3}{4}{5}{2}.
\item Find $x$ if \rvabBox{x}{2}{3}{4} $=$ \rvabBox{1}{x}{5}{2}.
\item Find all integers $x$ such that \rvabBox{x}{x}{x}{x} $=50$.
\end{enumerate}
\end{probbox}
\probsol{
\textbf{Understanding the problem.} The rule: multiply the two numbers on each diagonal, then add the two products. Part a is an evaluation (with a negative number); parts b and c are equations.
\step{1}{Part a}
\rvabBox{-3}{4}{5}{2} $=(-3)\cdot 2+4\cdot 5=-6+20=14$.
\begin{recallbox}The product of a negative number and a positive number is negative: $(-3)\cdot 2=-6$. Adding numbers with different signs: $-6+20=20-6=14$.\end{recallbox}
\step{2}{Part b: write both sides using the rule}
Left side: $x\cdot 4+2\cdot 3=4x+6$.\quad Right side: $1\cdot 2+x\cdot 5=2+5x$.\quad Equation: $4x+6=5x+2$.
\step{3}{Solve part b}
Subtract $4x$ from both sides: $6=x+2$. Subtract $2$ from both sides: $x=4$.
\verify{Left side $4\cdot 4+6=22$; right side $2+5\cdot 4=22$. Equal.}
\step{4}{Part c}
\rvabBox{x}{x}{x}{x} $=x\cdot x+x\cdot x=2x^2$. The equation is $2x^2=50$, that is, $x^2=25$.
Which numbers squared give $25$? There is $5\cdot 5=25$, but also $(-5)\cdot(-5)=25$. So $x=5$ or $x=-5$.
\begin{warnbox}Many students find only $x=5$ and forget $x=-5$. The product of two negative numbers is positive, so $(-5)^2=25$. The question asks for ``all integers'', which hints that there may be more than one answer.\end{warnbox}
\verify{$x=-5$: $(-5)(-5)+(-5)(-5)=25+25=50$. Correct.}
\ans{a) $14$;\ \ b) $x=4$;\ \ c) $x=5$ or $x=-5$.}
}

% ---------------------------------------------------------------- Problem 34
\begin{probbox}{34}{\stars{1}}
\begin{enumerate}[label=\alph*),leftmargin=1.6em,itemsep=1pt,topsep=0pt]
\item The sum of five consecutive odd numbers is $155$. Find the five numbers.
\item Leo says: ``I have found four consecutive odd numbers whose sum is $150$.'' Is Leo right or wrong? Why?
\end{enumerate}
\end{probbox}
\probsol{
\textbf{Understanding the problem.} Consecutive odd numbers differ by $2$ (for example $29$ and $31$). Part a asks us to find numbers; part b asks for an \emph{explanation}, not just a guess.
\step{1}{Part a: let the middle number be the unknown}
Call the middle number $n$. The five numbers are $n-4,\ n-2,\ n,\ n+2,\ n+4$.
\step{2}{Write and solve the equation}
$(n-4)+(n-2)+n+(n+2)+(n+4)=155$. The $-4$ and $+4$, and the $-2$ and $+2$, cancel: $5n=155$, so $n=31$.\\
The five numbers are $27,\ 29,\ 31,\ 33,\ 35$.
\verify{$27+29+31+33+35=155$. Correct, and they are consecutive odd numbers.}
\step{3}{Part b: try to write an equation}
Let the smallest odd number be $m$. The four numbers are $m,\ m+2,\ m+4,\ m+6$. Sum: $4m+12=150$, so $4m=138$, $m=138\div 4=34.5$.
The number $34.5$ is not an integer, let alone an odd number. So there are \emph{no} four consecutive odd numbers with sum $150$.
\begin{proofbox}Another explanation using even and odd: the sum of two odd numbers is even, so the sum of four odd numbers is even -- this is not yet a contradiction. But we can go further: $4m+12=4(m+3)$ is always \emph{divisible by $4$}, while $150$ is not divisible by $4$ ($150=4\cdot 37+2$). So the sum of four consecutive odd numbers can never be $150$.\end{proofbox}
\begin{warnbox}Trying a few sets of numbers (such as $35+37+39+41=152$) and then concluding is not rigorous enough. We need a reason that works in \emph{every} case.\end{warnbox}
\ans{a) $27,\ 29,\ 31,\ 33,\ 35$;\ \ b) Leo is wrong, because the sum of four consecutive odd numbers is always divisible by $4$, but $150$ is not.}
}

% ---------------------------------------------------------------- Problem 35
\begin{probbox}{35}{\stars{1}}
Every $2$ scoops of fish food are enough to feed $12$ fish for one day.
\begin{enumerate}[label=\alph*),leftmargin=1.6em,itemsep=1pt,topsep=2pt]
\item How many fish can $5$ scoops of food feed for one day?
\item Mia's fish tank has $18$ fish. A box of food contains $60$ scoops. For how many days will the box feed the whole tank?
\end{enumerate}
\end{probbox}
\probsol{
\textbf{Understanding the problem.} For one day, more scoops feed more fish, and doubling the scoops doubles the fish: the number of fish fed for one day is \emph{directly proportional} to the number of scoops.
\begin{recallbox}Two quantities $x$ and $y$ are directly proportional when $y=k\cdot x$ for some constant $k\neq 0$. Then the ratio $\dfrac{y}{x}$ is always $k$, so we can write the proportion $\dfrac{y_1}{x_1}=\dfrac{y_2}{x_2}$.\end{recallbox}
\step{1}{Find the constant of proportionality}
Let $y$ be the number of fish fed for one day and $x$ the number of scoops. When $x=2$, $y=12$, so $k=\dfrac{12}{2}=6$, that is $y=6x$: each scoop feeds $6$ fish for one day.
\step{2}{Part a}
For $x=5$: $y=6\cdot 5=30$ (fish).\\
(Or write the proportion $\dfrac{12}{2}=\dfrac{y}{5}$, so $y=\dfrac{12\cdot 5}{2}=30$.)
\step{3}{Part b: how many scoops does the tank need each day?}
Let $z$ be the number of scoops needed for $18$ fish for one day. Then $6z=18$, so $z=18\div 6=3$ (scoops).
\begin{center}
\begin{tabular}{|l|c|c|c|}
\hline
\rowcolor{tmBlue!10} Scoops $x$ & $2$ & $5$ & $z=3$\\ \hline
Fish $y=6x$ & $12$ & $30$ & $18$\\ \hline
\end{tabular}
\end{center}
\step{4}{Number of days}
The box holds $60$ scoops and $3$ scoops are used each day. Let the number of days be $d$: $3d=60$, so $d=60\div 3=20$ (days).
\begin{warnbox}A common mistake is to calculate $60\div 18$ or $60\cdot 6\div 18$ without understanding what it means. Always ask yourself: ``How many scoops are used each day?'' and only then divide.\end{warnbox}
\verify{$20$ days $\times$ $3$ scoops $=60$ scoops, exactly one box; and $3$ scoops $\times 6=18$ fish each day.}
\begin{parentbox}The unit ``fish-day'' can help: each scoop provides $6$ ``fish-days''; the whole box provides $60\cdot 6=360$ fish-days; dividing by $18$ fish gives $20$ days.\end{parentbox}
\ans{a) $30$ fish;\ \ b) $20$ days.}
}

% ---------------------------------------------------------------- Problem 36
\begin{probbox}{36}{\relax}
Each bus ride costs Lucy \$2.50. The bus company also sells a monthly pass for \$72 with unlimited rides.
\begin{enumerate}[label=\alph*),leftmargin=1.6em,itemsep=1pt,topsep=2pt]
\item Write an expression for the amount (in dollars) Lucy pays if she makes $n$ rides in a month without a monthly pass.
\item What is the smallest number of rides in a month for which the monthly pass is cheaper than paying per ride?
\end{enumerate}
\end{probbox}
\probsol{
\textbf{Understanding the problem.} Paying per ride: the cost depends on the number of rides $n$. Monthly pass: always \$72. We want the smallest number of rides for which paying per ride is \emph{more expensive} than the pass.
\step{1}{Part a}
Each ride costs \$2.50, so $n$ rides cost $2.5\cdot n=2.5n$ (dollars).
\step{2}{Part b: write an inequality}
The pass is cheaper when the cost of single rides is more than \$72, that is $2.5n>72$. Divide both sides by the positive number $2.5$: $n>72\div 2.5=28.8$.
\step{3}{Choose the smallest whole number and check}
The smallest whole number greater than $28.8$ is $29$. Check with a table:
\begin{center}
\begin{tabular}{|c|c|c|}
\hline
\rowcolor{tmBlue!10} Rides $n$ & Per ride $2.5n$ (dollars) & Compared with the \$72 pass\\ \hline
$28$ & $70$ & paying per ride is cheaper\\ \hline
$29$ & $72.5$ & \textbf{the pass is cheaper}\\ \hline
\end{tabular}
\end{center}
So she needs at least $29$ rides.
\begin{warnbox}The number of rides is a whole number, so do not answer ``$28.8$ rides''. Do not ``round'' to $28$ either: $28$ rides cost only \$70, which is still cheaper than the pass. We need the smallest whole number \emph{greater than} $28.8$, which is $29$.\end{warnbox}
\begin{parentbox}This is a great ``real life'' problem to talk about with your child: phone plans, monthly swimming passes... can all be compared in this way.\end{parentbox}
\ans{a) $2.5n$ dollars;\ \ b) at least $29$ rides.}
}

% ---------------------------------------------------------------- Problem 37
\begin{probbox}{37}{\relax}
Look at the table below (the values increase steadily as $n$ increases):
\begin{center}
\begin{tabular}{|c|c|c|c|c|c|}
\hline
\rowcolor{tmBlue!10} $n$ & $1$ & $2$ & $3$ & $4$ & $5$\\ \hline
Value & $2$ & $7$ & $12$ & $17$ & $22$\\ \hline
\end{tabular}
\end{center}
\begin{enumerate}[label=\alph*),leftmargin=1.6em,itemsep=1pt,topsep=2pt]
\item Write an expression in $n$ for the row of values.
\item Find the value for $n=40$.
\item For which $n$ is the value $302$?
\item Is any value in the table (continued forever) equal to $2026$?
\end{enumerate}
\end{probbox}
\probsol{
\textbf{Understanding the problem.} Each time $n$ increases by $1$, the value increases by $5$. It is an arithmetic sequence with first term $2$ and common difference $5$.
\begin{recallbox}If the value always increases by $d$ when $n$ increases by $1$, the expression has the form $dn+b$; the $n$-th term of an arithmetic sequence is $a+(n-1)d$.\end{recallbox}
\step{1}{Part a}
The expression has the form $5n+b$. Substitute $n=1$: $5+b=2$, so $b=-3$. Expression: $5n-3$.
(Or: $2+5(n-1)=2+5n-5=5n-3$.)
\verify{$n=4$: $5\cdot 4-3=17$; \ $n=5$: $5\cdot 5-3=22$. It matches the table.}
\step{2}{Part b}
$5\cdot 40-3=200-3=197$.
\step{3}{Part c: solve an equation}
$5n-3=302\ \Rightarrow\ 5n=305\ \Rightarrow\ n=61$.
\step{4}{Part d}
$5n-3=2026\ \Rightarrow\ 5n=2029$. But $2029$ ends in $9$, so it is not divisible by $5$, and there is no whole number $n$ that works.
\begin{proofbox}Every value $5n-3$ ends in $2$ or $7$ (because $5n$ ends in $0$ or $5$). The number $2026$ ends in $6$, so it cannot be in the table.\end{proofbox}
\ans{a) $5n-3$;\ \ b) $197$;\ \ c) $n=61$;\ \ d) no.}
}

% ---------------------------------------------------------------- Problem 38
\begin{probbox}{38}{\stars{1}}
\begin{enumerate}[label=\alph*),leftmargin=1.6em,itemsep=1pt,topsep=0pt]
\item The product $60\times 24\times 365$ is the number of which time unit in a non-leap year (with $365$ days)? Work out the product.
\item How many days, hours, minutes and seconds are there in one million seconds?
\end{enumerate}
\end{probbox}
\probsol{
\textbf{Understanding the problem.} Part a: read the product as a chain of unit conversions. Part b: convert back from the small unit (seconds) to larger units, using division with remainder.
\begin{recallbox}$1$ minute $=60$ seconds; $1$ hour $=60$ minutes $=3\,600$ seconds; $1$ day $=24$ hours $=24\cdot 3\,600=86\,400$ seconds.\end{recallbox}
\step{1}{Part a}
$365$ is the number of days in the year; $24\times 365$ is the number of hours in the year; $60\times 24\times 365$ is the number of \textbf{minutes} in a non-leap year.
Calculation: $24\times 365=8\,760$; \ $60\times 8\,760=525\,600$ (minutes).
\step{2}{Part b: convert to days}
$1\,000\,000\div 86\,400$ gives quotient $11$, remainder $1\,000\,000-11\cdot 86\,400=1\,000\,000-950\,400=49\,600$ (seconds).
\step{3}{Convert the remainder to hours}
$49\,600\div 3\,600$ gives quotient $13$, remainder $49\,600-13\cdot 3\,600=49\,600-46\,800=2\,800$ (seconds).
\step{4}{Convert the remainder to minutes and seconds}
$2\,800\div 60$ gives quotient $46$, remainder $2\,800-46\cdot 60=2\,800-2\,760=40$ (seconds).
\begin{center}
\begin{tabular}{|l|r|r|}
\hline
\rowcolor{tmBlue!10} Division & Quotient & Remainder (seconds)\\ \hline
$1\,000\,000\div 86\,400$ & $11$ days & $49\,600$\\ \hline
$49\,600\div 3\,600$ & $13$ hours & $2\,800$\\ \hline
$2\,800\div 60$ & $46$ minutes & $40$\\ \hline
\end{tabular}
\end{center}
\verify{$11\cdot 86\,400+13\cdot 3\,600+46\cdot 60+40=950\,400+46\,800+2\,760+40=1\,000\,000$.}
\begin{warnbox}Do not write the result as a decimal such as ``$11.57$ days'' and then read it as ``$11$ days $57$ hours'': the decimal part of a day is not a number of hours. Use division with remainder as above.\end{warnbox}
\ans{a) The number of minutes in a non-leap year, $525\,600$;\ \ b) $11$ days $13$ hours $46$ minutes $40$ seconds.}
}

% ---------------------------------------------------------------- Problem 39
\begin{probbox}{39}{\stars{1}}
Four positive integers $a$, $b$, $c$, $d$ satisfy
\[\frac{3}{a-5}=\frac{3}{b+1}=\frac{3}{c-2}=\frac{3}{d+4}\qquad\text{and}\qquad a+b+c+d=50.\]
\begin{enumerate}[label=\alph*),leftmargin=1.6em,itemsep=1pt,topsep=2pt]
\item Arrange $a$, $b$, $c$, $d$ in increasing order.
\item Find the value of each number.
\end{enumerate}
\end{probbox}
\probsol{
\textbf{Understanding the problem.} Four fractions with the same numerator $3$ are equal. Use the ``set it equal to $k$'' technique as in Problem 12, then use the condition on the sum to find $k$.
\step{1}{The denominators are equal}
\begin{recallbox}If $\dfrac{m}{p}=\dfrac{m}{q}$ with $m\neq 0$, then $p=q$.\end{recallbox}
$a-5=b+1=c-2=d+4$. Call the common value $k$ ($k\neq 0$ because a denominator cannot be $0$).
\step{2}{Write each number in terms of $k$}
$a=k+5$,\quad $b=k-1$,\quad $c=k+2$,\quad $d=k-4$.
\step{3}{Part a: compare}
$k-4<k-1<k+2<k+5$, so $d<b<c<a$.
\step{4}{Part b: use the sum to find $k$}
$a+b+c+d=(k+5)+(k-1)+(k+2)+(k-4)=4k+2$. From $4k+2=50$ we get $4k=48$, $k=12$.\\
So $a=17$, \ $b=11$, \ $c=14$, \ $d=8$.
\verify{$a-5=12$, $b+1=12$, $c-2=12$, $d+4=12$: all four fractions equal $\dfrac{3}{12}$. The sum is $17+11+14+8=50$. And $8<11<14<17$.}
\begin{warnbox}When adding $(+5)+(-1)+(+2)+(-4)$, keep the signs: the result is $+2$, not $12$.\end{warnbox}
\ans{a) $d<b<c<a$;\ \ b) $a=17$, $b=11$, $c=14$, $d=8$.}
}

% ---------------------------------------------------------------- Problem 40
\begin{probbox}{40}{\stars{2}}
There are four bags of rice with different weights. They are weighed two at a time in all possible ways, giving six results:
\[7 \text{ kg},\quad 8 \text{ kg},\quad 9 \text{ kg},\quad 10 \text{ kg},\quad 11 \text{ kg},\quad 12 \text{ kg}.\]
\begin{enumerate}[label=\alph*),leftmargin=1.6em,itemsep=1pt,topsep=2pt]
\item Find the total weight of the four bags.
\item Each bag weighs a whole number of kilograms. Find the weight of each bag.
\end{enumerate}
\end{probbox}
\probsol{
\textbf{Understanding the problem.} Let the weights of the four bags from lightest to heaviest be $a<b<c<d$ (kg). Four bags can be paired in $6$ ways:
$a+b,\ a+c,\ a+d,\ b+c,\ b+d,\ c+d$. This problem extends Problem 15 (the three pumpkins).
\step{1}{Part a: add all six results}
Each bag appears in exactly $3$ pairs (paired with each of the $3$ other bags). Therefore
$3(a+b+c+d)=7+8+9+10+11+12=57$, \ so $a+b+c+d=57\div 3=19$ (kg).
\begin{recallbox}With three objects, each object appears in $2$ pairs, so the total $=$ (sum of the pairs)$\div 2$. With four objects, each appears in $3$ pairs, so the total $=$ (sum of the pairs)$\div 3$.\end{recallbox}
\step{2}{Which pairs certainly match which results?}
\begin{itemize}[leftmargin=1.4em,itemsep=1pt,topsep=1pt]
\item The lightest pair is the two lightest bags: $a+b=7$.
\item The second lightest pair is $a+c=8$: every other pair (except $a+b$) is heavier than $a+c$, because $a+d>a+c$, \ $b+c>a+c$, \ $b+d>a+c$, \ $c+d>a+c$.
\item Similarly, the heaviest pair is $c+d=12$, and the second heaviest is $b+d=11$.
\item The two middle results ($9$ and $10$) are $a+d$ and $b+c$, but we do \emph{not yet know} which is $9$ and which is $10$.
\end{itemize}
\begin{proofbox}Subtract $a+b=7$ from $a+c=8$: $c-b=1$, that is, $c=b+1$. Then $b+c=b+(b+1)=2b+1$ is an \emph{odd number} (because $b$ is a whole number). Of the two numbers $9$ and $10$, only $9$ is odd. So $b+c=9$ and $a+d=10$.\end{proofbox}
\step{3}{Find each bag}
$2b+1=9\Rightarrow b=4$;\quad $c=b+1=5$;\quad $a=7-b=3$;\quad $d=12-c=7$.
\begin{center}
\begin{tabular}{|c|c|c|c|c|c|c|}
\hline
\rowcolor{tmBlue!10} Pair & $a+b$ & $a+c$ & $b+c$ & $a+d$ & $b+d$ & $c+d$\\ \hline
Sum (kg) & $3+4=7$ & $3+5=8$ & $4+5=9$ & $3+7=10$ & $4+7=11$ & $5+7=12$\\ \hline
\end{tabular}
\end{center}
\verify{$3+4=7$, $3+5=8$, $3+7=10$, $4+5=9$, $4+7=11$, $5+7=12$: exactly the six results $7,8,9,10,11,12$. Total $3+4+5+7=19$.}
\begin{warnbox}Do not simply assume that the third result ($9$ kg) is $a+d$ or $b+c$. You need a reason! Without the ``whole number'' condition, the case $b+c=10$ giving $a=2.5$, $b=4.5$, $c=5.5$, $d=6.5$ would also work.\end{warnbox}
\begin{parentbox}You can experiment with four slips of paper showing $3,4,5,7$: let your child pair them in all possible ways and add, and they will see for themselves that each number is used three times.\end{parentbox}
\ans{a) $19$ kg;\ \ b) $3$ kg, $4$ kg, $5$ kg, $7$ kg.}
}
